% Generated draft: requires mathematical and source review.
\documentclass[11pt]{article}
\usepackage[margin=25mm]{geometry}
\usepackage{fontspec}
\setmainfont{Latin Modern Roman}
\usepackage{amsmath,amssymb,mathtools,unicode-math}
\setmathfont{Latin Modern Math}
\usepackage{enumitem,needspace,xurl,hyperref}
\setlength{\parindent}{0pt}
\setlength{\parskip}{5pt}
\setlength{\emergencystretch}{3em}
\newcommand{\N}{\mathbb{N}}
\newcommand{\Z}{\mathbb{Z}}
\newcommand{\Q}{\mathbb{Q}}
\newcommand{\R}{\mathbb{R}}
\newcommand{\C}{\mathbb{C}}
\newcommand{\F}{\mathbb{F}}
\newcommand{\A}{\mathbb{A}}
\newcommand{\PP}{\mathbb{P}}
\newcommand{\E}{\mathbb{E}}
\newcommand{\eps}{\varepsilon}
\newcommand{\dd}{\,\mathrm{d}}
\begin{document}
\section*{UnsolvedMath ID: 4900013}
\textbf{Source problem code:} catalogue entry
\subsection{Definitions and mathematical statement}
Let $P\subset\R^2$ be a compact convex polygon with nonempty interior. For a point $x$ outside $P$ in the regular domain, choose a supporting tangent line with a fixed orientation convention and tangency vertex $v$. The outer-billiard map sends $x$ to $T_P(x)=2v-x$. Points on singular supporting rays, where the tangency choice is not unique, must be excluded or assigned a specified convention.

A forward orbit is defined if $T_P^k(x)$ avoids these singularities for every integer $k\geq0$. It is unbounded if
\[
\sup_{k\geq0}\|T_P^k(x)\|_2=+\infty.
\]
The source conjecture says that for almost every convex polygon $P$ there exists such an unbounded orbit. This quantifier is existential in $x$; it does not assert that almost every orbit is unbounded.

A standard possible parameter convention fixes a number $n\geq3$ of vertices and uses Lebesgue measure on the open set of labelled, cyclically ordered strictly convex $n$-gons in $(\R^2)^n$. With that convention, the target for an admissible $n$ would be that the polygons possessing an unbounded nonsingular orbit form a full-measure subset. The source does not specify the intended values of $n$ or a measure on polygon space. There is no canonical single Lebesgue measure on polygons of all numbers of vertices, so this convention and any exceptional polygon classes must be settled before a precise almost-everywhere statement can be asserted. Whether ``orbit'' means forward or two-sided is also not specified; the display above makes the forward interpretation explicit.
\subsection{Short English statement}
Does a typical convex polygon admit at least one unbounded outer-billiard orbit? The number of sides, measure on polygon space, and singular-orbit conventions need to be specified.
\subsection{Sources}
\begin{itemize}
\item Local UnsolvedMath registry, permanent ID 4900013, exact source code catalogue entry, ``Unbounded outer-billiard orbits for almost every polygon''; no proposer or year supplied. Registry status partially solved, classification PARTIAL-PROGRESS.
\item Wikipedia contributors, \emph{List of unsolved problems in mathematics}, Dynamical systems, source bullet 13, \url{https://en.wikipedia.org/wiki/List_of_unsolved_problems_in_mathematics#Dynamical_systems}. Supplied HTML extraction accessed July 29, 2026. The list calls the item unsolved but is not an authoritative resolution source; provenance flags status NEEDS\_REVIEW. No primary citation for the exact measure-generic conjecture is supplied.
\item The supplied research summary reports unbounded orbits for irrational kites (Schwartz) and semidisks/near-semidisks (Dolgopyat--Fayyad), and boundedness of all orbits for quasi-rational polygons (Vivaldi--Shaidenko, Ko\l odziej, Gutkin--Simanyi). It describes the original Moser--Neumann existence question as answered positively while leaving the measure-generic polygon conjecture open. No titles, dates, hypotheses, or precise citations are supplied; semidisk examples are related nonpolygonal cases, not the claimed full-measure polygon result.
\item Rights note from the source: Wikipedia text is available under CC BY-SA 4.0; attribution and share-alike compatibility require release review. This draft does not resolve that release question.
\end{itemize}
\end{document}
